\documentclass[10pt,a4paper]{article}

\usepackage[T1]{fontenc}
\usepackage[utf8]{inputenc}
\usepackage{tgtermes}
\usepackage{tgheros}
\usepackage{microtype}
\usepackage[a4paper,top=18mm,bottom=20mm,left=18mm,right=18mm]{geometry}
\usepackage{xcolor}
\usepackage{graphicx}
\usepackage{booktabs}
\usepackage{tabularx}
\usepackage{array}
\usepackage{enumitem}
\usepackage[most]{tcolorbox}
\usepackage{fancyhdr}
\usepackage{lastpage}
\usepackage{titlesec}
\usepackage{hyperref}
\usepackage{bookmark}
\usepackage{ragged2e}
\usepackage{caption}
\usepackage{placeins}
\usepackage{amsmath}
\usepackage{amssymb}

\definecolor{Navy}{HTML}{17324D}
\definecolor{Teal}{HTML}{008C8C}
\definecolor{TealLight}{HTML}{E1F2F2}
\definecolor{Orange}{HTML}{E58B2A}
\definecolor{OrangeLight}{HTML}{FFF1DE}
\definecolor{Green}{HTML}{2E7D5B}
\definecolor{GreenLight}{HTML}{E8F4ED}
\definecolor{Red}{HTML}{B94A3D}
\definecolor{RedLight}{HTML}{F9E9E7}
\definecolor{Slate}{HTML}{607284}
\definecolor{Light}{HTML}{F2F6F8}
\definecolor{Line}{HTML}{D8E2E9}
\definecolor{Ink}{HTML}{172A3A}

\hypersetup{
  colorlinks=true,
  linkcolor=Teal,
  urlcolor=Teal,
  pdftitle={HeiQ XMembrane -- Model Closure and Material Handoff},
  pdfauthor={XMembrane project team | HeiQ Materials AG},
  pdfsubject={Closure of the broad XMembrane model phase and handoff to material and coupon work},
  pdfkeywords={XMembrane, electro-osmosis, ERV, model closure, membrane, material search, coupon}
}

\setlength{\parindent}{0pt}
\setlength{\parskip}{5pt}
\raggedbottom
\setlist[itemize]{leftmargin=5mm,itemsep=2.2pt,topsep=3pt}
\setlist[enumerate]{leftmargin=6mm,itemsep=2.3pt,topsep=3pt}
\newcolumntype{Y}{>{\RaggedRight\arraybackslash}X}
\newcolumntype{C}{>{\Centering\arraybackslash}X}
\renewcommand{\arraystretch}{1.16}
\captionsetup{font=small,labelfont={bf,sf,color=Navy}}

\titleformat{\section}
  {\sffamily\bfseries\color{Navy}\fontsize{17}{20}\selectfont}
  {\thesection}{8pt}{}
\titlespacing*{\section}{0pt}{9pt}{5pt}
\titleformat{\subsection}
  {\sffamily\bfseries\color{Navy}\fontsize{12.5}{15}\selectfont}
  {\thesubsection}{7pt}{}

\pagestyle{fancy}
\fancyhf{}
\fancyhead[L]{\sffamily\footnotesize\color{Slate}HEIQ MATERIALS AG · XMEMBRANE}
\fancyhead[R]{\sffamily\footnotesize\color{Slate}MODEL CLOSURE · 31.08.2026}
\fancyfoot[L]{\sffamily\scriptsize\color{Slate}Internal model decision basis · not product validation}
\fancyfoot[R]{\sffamily\scriptsize\color{Slate}Page \thepage\ of \pageref*{LastPage}}
\renewcommand{\headrulewidth}{0.35pt}

\newtcolorbox{mainbox}[1][]{
  enhanced,colback=TealLight,colframe=Teal,boxrule=0.8pt,arc=1.5mm,
  left=4mm,right=4mm,top=3mm,bottom=3mm,#1
}
\newtcolorbox{warningbox}[1][]{
  enhanced,colback=OrangeLight,colframe=Orange,boxrule=0.8pt,arc=1.5mm,
  left=4mm,right=4mm,top=3mm,bottom=3mm,#1
}
\newtcolorbox{stopbox}[1][]{
  enhanced,colback=RedLight,colframe=Red,boxrule=0.8pt,arc=1.5mm,
  left=4mm,right=4mm,top=3mm,bottom=3mm,#1
}
\newtcolorbox{plainbox}[1][]{
  enhanced,colback=Light,colframe=Line,boxrule=0.6pt,arc=1.5mm,
  left=4mm,right=4mm,top=3mm,bottom=3mm,#1
}
\newcommand{\eyebrow}[1]{{\sffamily\bfseries\small\color{Teal}\MakeUppercase{#1}}}
\newcommand{\evidencepath}[1]{\texttt{\small\detokenize{#1}}}

\begin{document}

% -----------------------------------------------------------------------------
% Cover
% -----------------------------------------------------------------------------
\thispagestyle{empty}
\pdfbookmark[0]{Cover}{cover}
\vspace*{-10mm}
\begin{tcolorbox}[enhanced,colback=Navy,colframe=Navy,boxrule=0pt,
  arc=0mm,left=11mm,right=11mm,top=13mm,bottom=13mm,width=\textwidth]
  {\sffamily\bfseries\color{TealLight}\large HEIQ XMEMBRANE PROJECT}\par\vspace{7mm}
  {\raggedright\sffamily\bfseries\color{white}\fontsize{29}{33}\selectfont
  Model closure and material handoff\par}
  \vspace{4mm}
  {\sffamily\color{white!85}\fontsize{14}{18}\selectfont
  Frozen energy basis · direct-ERV gradient closure · AC null test · next experimental gate}
\end{tcolorbox}

\vspace{13mm}
{\sffamily\bfseries\fontsize{16}{20}\selectfont\color{Navy}Internal technical decision note · 31 August 2026}\par
\vspace{4mm}
{\fontsize{12}{17}\selectfont\color{Ink}
This is the current authoritative interpretation of the model study. It
supersedes earlier reports where a positive passive surface-gradient direction
was treated as a feasibility requirement or water diffusivity was presented as
an independently sufficient material criterion. Earlier reports remain the
chronological record.}

\vspace{8mm}
\begin{mainbox}
\textbf{Phase decision.} The broad steady-state model phase is complete.
Start the joint-property material search and controlled-flow coupon work now.
Additional unfocused parameter sweeps are unlikely to change that decision.
The project has identified conditional opportunity and the measurements needed
to test it; it has not validated a membrane, selected a product architecture or
demonstrated a net system energy saving.
\end{mainbox}

\vfill
\begin{tabularx}{\textwidth}{@{}p{43mm}Y@{}}
\textbf{Primary comparator} & Passive ERV + remaining AC = 543.741~kW; cooling is included \\
\textbf{Main material result} & Performance is governed by a coupled property bundle, not by $D$ or transport number alone \\
\textbf{Direct-ERV result} & Active uphill transport is thermodynamically admissible; film losses and the counter-current pinch determine the local gradient regime \\
\textbf{Symmetric-current result} & No useful quasi-static mean EO transport beyond the passive $N(0)$ state in the present model \\
\textbf{Current/area boundary} & No current-density acceptance ceiling and no exploratory membrane-area acceptance gate
\end{tabularx}

\clearpage

% -----------------------------------------------------------------------------
\pdfbookmark[0]{Abstract}{abstract}
\section*{Abstract}
\addcontentsline{toc}{section}{Abstract}

The frozen O80 equal-service basis is 1,166.400~kW for AC alone and
543.741~kW for passive ERV plus remaining AC. The post-ERV O30
single-common-voltage coordinate gives a 381.327-kW known subtotal and a
162.414-kW conditional window, while the full-duty direct-ERV coordinate gives
107.856~kW and a 435.884-kW (80.16\%) conditional window. These differences
exclude one or more of receiver/sink work, EO heat rejection, electrodes,
manifolds, leakage, pressure drop and durability; they are not net savings.

The final direct-ERV audit resolves the gradient question. Positive passive
surface direction is a regime diagnostic, not a thermodynamic gate for active
EO. At the frozen O30 mechanism coordinate, the bulk activity margin is
favourable, but two finite gas films reverse the surface margin to
$-0.00797$. A 1.137-fold common increase in film capacity, or the equivalent
diagnostic area increase at fixed duty, restores a non-negative surface
direction. At full duty the bulk counter-current pinch is exactly zero, so no
finite film or area multiplier produces a positive surface margin at finite
flux; EO must act uphill at the pinch. All signed states close the applicable
thermodynamic audit.

A matched quasi-static $+J/0/-J$ replay also closes the proposed symmetric-AC
question within the current constant-drag model. O30 produces only
0.00152~kg/h of even-in-current transport above the passive state, equal to
0.348~ppm of the $+J$ water rate, while full duty gives zero. Thus a
charge-balanced symmetric current does not provide useful mean EO pumping here;
transient or rectified AC would require measured non-linearity, asymmetry or
memory not present in this model. The correct next step is same-sample
measurement of sorption, chemical water permeance, drag, wet conductivity,
compressed-stack ASR, thickness, film exchange and coupled water flux in a
controlled coupon. Candidate ranking and full-system claims remain open.

\begin{mainbox}
\textbf{Closure statement.} We know enough to stop the broad modelling loop.
The remaining uncertainties are empirical and subsystem-specific; they should
be resolved with materials, coupons and a receiver/thermal ledger rather than
by searching a larger hypothetical parameter space.
\end{mainbox}

% -----------------------------------------------------------------------------
\section{Frozen energy and comparison basis}
\eyebrow{Use equal service and state every omitted load}

\begin{center}\small
\begin{tabularx}{\textwidth}{@{}Y r r Y@{}}
\toprule
\textbf{System coordinate} & \textbf{Power} & \textbf{Difference vs ERV+AC} & \textbf{Claim class} \\
 & [kW] & [kW / \%] & \\
\midrule
AC only, O80 & 1,166.400 & $-622.659$ / $-114.51$ & Calibrated legacy comparator \\
Passive ERV + AC & 543.741 & 0 / 0 & Primary comparator: 34.359-kW ERV + 509.381-kW cooling \\
Post-ERV EO to O30, common 1.00~V & 381.327 & $+162.414$ / $+29.87$ & Conditional incomplete ledger \\
Post-ERV EO to O30, uniform-current cross-check & 379.516 & $+164.225$ / $+30.20$ & Conditional topology cross-check \\
Direct EO inside ERV, full humidity duty, 0.25~V & 107.856 & $+435.884$ / $+80.16$ & Conditional architecture screen \\
\bottomrule
\end{tabularx}
\end{center}

For the post-ERV common-voltage row, 316.910~kW is remaining cooling and
64.417~kW is the ERV/active-device subtotal. For the direct-ERV row,
44.438~kW is remaining cooling and 63.419~kW is the known active-core/fan
subtotal. The conservative post-ERV all-in break-even is
0.1567008~kWh/kg water removed (156.7~Wh/kg), not the historical
180-Wh/L shorthand.

\begin{warningbox}
The positive numbers in the difference column are \textbf{available conditional
windows}, not measured or net system savings. The direct-ERV 80.16\% result is
the largest upside and the least complete system boundary. It assumes
prescribed 0.75 sensible effectiveness, external EO-heat rejection, idealised
hydraulics and an unmeasured inverse-design material vector.
\end{warningbox}

The final common-voltage replay below 2~V separates numerical status from
physical interpretation. On the refined 16-segment grid, O30 has 14 converged
cases, 14 physical no-solution coordinates and 21 numerical/local unresolved
coordinates; all seven sampled diffusivities nevertheless have a confirmed
selected O30 case at 1.00--1.25~V. Full remaining duty has zero confirmed
cases, 21 physical no-solution coordinates and 28 numerical/local unresolved
coordinates. Across both segment counts, all 98 cases that encountered the
legacy high-bracket flag were retried after increasing the configured numerical
current bound from 50,000 to 5,000,000~A/m$^2$; none recovered. This rules out
the configured numerical bound as the active limiter in those retries, but it
does not convert the unresolved rows into a global physical no-go proof.

The post-ERV O30 known-core coordinate removes 1,228.27~kg/h at
30.057~kW (40.864~kg/kWh known-core FOM). The direct-ERV full-duty coordinate
transfers 4,827.15~kg/h with a 63.419-kW known active-core/fan subtotal. These
FOMs have different device boundaries and must not be used as product
efficiencies without the omitted loads.

% -----------------------------------------------------------------------------
\section{Direct EO in the ERV: corrected gradient interpretation}
\eyebrow{Active feasibility and passive assistance are different questions}

For net molar water flux $N>0$ from source to receiver, the finite-film model
obeys the segment identity
\begin{equation}
  \Delta a_{\mathrm{surface}}
  = \Delta a_{\mathrm{bulk}}
  - N\left(\frac{1}{K_s}+\frac{1}{K_r}\right),
  \label{eq:film-identity}
\end{equation}
where $K_s$ and $K_r$ are the source- and receiver-side film capacities.
The numerical identity error is below $3.0\times10^{-15}$ in the closure
replay. A negative $\Delta a_{\mathrm{surface}}$ says that EO is transporting
uphill locally; it is not by itself a violation of thermodynamics. The active
state is admissible when the electrical/chemical coupling and dissipation
audit closes.

\begin{center}\small
\begin{tabularx}{\textwidth}{@{}p{34mm}r r r r Y@{}}
\toprule
\textbf{Coordinate} & \textbf{Bulk $\Delta a$} & \textbf{Surface $\Delta a$} & \textbf{Area} & \textbf{$V_t$} & \textbf{Interpretation} \\
\midrule
Direct-ERV O30 mechanism point & $+0.0581$ to $+0.0969$ & $-0.00797$ to $-0.00536$ & 24,604~m$^2$ & 0.025~V & Passive bulk assistance exists, but film losses reverse the local membrane-surface direction \\
Direct-ERV full duty & 0 at the pinch & $-0.0822$ to $-0.0495$ & 34,425~m$^2$ & 0.25~V & No passive bulk headroom at the pinch; finite-flux transfer is locally uphill \\
\bottomrule
\end{tabularx}
\end{center}

At O30, multiplying both film capacities by 1.137 restores
$\min(\Delta a_{\mathrm{surface}})=0$. The equivalent area diagnostic at fixed
duty is 27,974~m$^2$. Recovering the former $+0.02$ robustness margin requires
1.733 times the film capacity or 42,641~m$^2$. These are diagnostic crossover
values, not area targets. At full duty, increasing film capacity or area only
approaches zero from below; no finite multiplier creates a positive surface
margin because the bulk margin is zero.

\begin{figure}[htbp]
  \centering
  \includegraphics[width=0.96\textwidth]{../../outputs/direct_erv_gradient_closure_v1/figures/01_film_area_gradient_envelope.pdf}
  \caption{Direct-ERV film/area gradient envelope. The O30 mechanism point has
  a finite surface-direction crossover; the full-duty counter-current pinch
  has none at finite local flux. Area is diagnostic, not an acceptance gate.}
\end{figure}
\FloatBarrier

Water-polarisation work accounts for up to 99.877\% of the 0.025-V O30
terminal voltage and 99.990\% of the 0.25-V full-duty voltage. This is why
``low resistance'' alone cannot identify a good membrane: the chemical water
work, passive permeance and interface conditions must be evaluated together.

\begin{plainbox}
\textbf{Resolved gradient statement.} Direct EO in the ERV can exploit a
favourable bulk gradient where one exists and can actively cross an adverse
surface gradient when electrical work closes the dissipation balance. The
potential benefit is therefore real at model level, but the film coefficients
and pinch location are first-order experimental inputs.
\end{plainbox}

% -----------------------------------------------------------------------------
\section{Symmetric-current AC: quasi-static closure}
\eyebrow{A gradient can break symmetry, but this model contains no useful rectifier}

The final test evaluates the same local constitutive law at $+J$, zero current
and $-J$, then averages the two signed-current states. The current amplitudes
are those of the frozen positive-DC coordinates. This is a quasi-static,
current-symmetric replay; it is not a frequency-domain or electrode model.

\begin{center}\small
\begin{tabularx}{\textwidth}{@{}p{31mm}r r r r r@{}}
\toprule
\textbf{Coordinate} & \textbf{$+J$ water} & \textbf{$N(0)$} & \textbf{Cycle mean} & \textbf{Increment over $N(0)$} & \textbf{Cycle electrical power} \\
 & [kg/h] & [kg/h] & [kg/h] & [kg/h] & [kW] \\
\midrule
Direct-ERV O30 & 4,365.918 & 0.415151 & 0.416669 & 0.00151783 & 20.092 \\
Direct-ERV full duty & 4,827.147 & 0 & 0 & 0 & 17.944 \\
\bottomrule
\end{tabularx}
\end{center}

The O30 even-in-current increment is only $3.48\times10^{-7}$ of the $+J$
water rate (0.348~ppm); full duty is exactly zero in the present model.
Consequently, a charge-balanced symmetric current adds no useful mean EO
transport beyond passive $N(0)$ while still dissipating electrical power.

\begin{figure}[htbp]
  \centering
  \includegraphics[width=0.98\textwidth]{../../outputs/direct_erv_gradient_closure_v1/figures/02_signed_triplet_ac_null.pdf}
  \caption{Matched $+J/0/-J$ quasi-static replay. Positive DC transports the
  required water, whereas signed-current averaging cancels the polarity-odd EO
  contribution and leaves essentially the passive state.}
\end{figure}
\FloatBarrier

This result does not exclude AC transport based on time-dependent hydration,
hysteresis, asymmetric interfaces, nonlinear electrode kinetics, induced
charge or a deliberately graded membrane. It says those mechanisms are absent
from the current evidence and should not trigger another modelling campaign
before dynamic measurements exist. Also, zero-mean voltage need not imply
zero-mean current when a chemical-potential offset is present.

% -----------------------------------------------------------------------------
\section{What the material search should optimise}
\eyebrow{The identifiable quantity is a coupled property bundle}

The focused $D=10^{-14}$--$10^{-9}$~m$^2$/s sweep held water capacity,
thickness, drag, conductivity, ASR and activity exponents fixed. It therefore
proves a strong \emph{one-factor sensitivity at that coordinate}; it does not
prove that a dry-state scalar diffusivity is universally dominant. In a
sorption-based membrane, the passive term contains the chemical water
permeance, schematically
\begin{equation}
  \mathcal{P}_{w,\mathrm{chem}}
  \propto \frac{D(a,T)}{L}
  \frac{\partial C_w}{\partial a_w}
  \quad\text{or an equivalent chemical-potential mobility.}
  \label{eq:chemical-permeance}
\end{equation}
The derivative, thickness and interfaces are therefore inseparable from $D$
when ranking candidates.

\begin{center}\small
\begin{tabularx}{\textwidth}{@{}p{45mm}Y Y@{}}
\toprule
\textbf{Same-sample property} & \textbf{Why it is required} & \textbf{Exploratory direction / method} \\
\midrule
Sorption $C_w(a,T)$, hysteresis and $\partial C_w/\partial\mu_w$ & Converts humidity into chemical loading/work and closes the passive mobility & Full adsorption/desorption curves across roughly $a=0.05$--0.9 and 23--32$^\circ$C \\
$D(a,T)$ or Maxwell--Stefan mobility, with definition & Separates kinetics from capacity and quantifies back-/co-diffusion & The post-ERV frozen screen flags $D\lesssim10^{-11}$ and prefers $\lesssim10^{-12}$~m$^2$/s; this is not a universal gate \\
EO drag $n_d(a,T,j)$ & Sets polarity-odd water transported per charge & High and reproducible; the inverse-design coordinate is about 100~mol H$_2$O/mol charge, not a minimum pass/fail value \\
Wet through-plane conductivity and compressed-stack ASR & Sets ohmic work and exposes contact losses & Same hydration/ion form as drag and diffusion; inverse-design coordinates are 23.17~S/m and $10^{-5}$~$\Omega\,\mathrm{m}^2$ \\
Dry/wet thickness, swelling, creep and strength & Couples permeance and resistance to durability & Thin but stable; inverse-design dry thickness is 36.9~$\mu$m \\
Source/receiver exchange coefficients & Determines whether bulk assistance survives at membrane surfaces & Measure each side versus humidity, flow and gap; direct-ERV O30 has a 1.137-fold crossover \\
Coupled net water flux and charge & Tests whether the property bundle coexists & Two-sided water balance under $+J/0/-J$ and repeated polarity blocks \\
\bottomrule
\end{tabularx}
\end{center}

Post-ERV polishing and direct-ERV integration need not prefer the same
material. Post-ERV EO generally benefits from suppressing adverse chemical
permeance. Direct ERV can use favourable passive transfer before the pinch, so
indiscriminately minimising permeance can remove useful free transport. Every
candidate should be replayed against both architecture boundaries after its
same-sample data are available.

\begin{warningbox}
Do not construct a synthetic ``best membrane'' by combining $D$, drag,
conductivity and thickness from different materials or conditioning states.
Nafion is not the target, area is not yet capped, there is no 21-A/m$^2$ gate,
and a transport number above 100 is not sufficient by itself.
\end{warningbox}

% -----------------------------------------------------------------------------
\section{Closed questions, open gates and handoff}
\eyebrow{Stop broad modelling; convert uncertainty into measurements}

\begin{center}\small
\begin{tabularx}{\textwidth}{@{}p{45mm}p{20mm}Y@{}}
\toprule
\textbf{Question or gate} & \textbf{Status} & \textbf{Decision} \\
\midrule
Energy comparator and cooling inclusion & \textcolor{Green}{\textbf{CLOSED}} & Use 543.741~kW passive ERV+AC as the primary equal-service basis \\
Active transport with a negative surface margin & \textcolor{Green}{\textbf{CLOSED}} & Thermodynamically admissible when the coupling/dissipation audit closes; passive direction remains a diagnostic \\
Direct-ERV film and pinch behaviour & \textcolor{Green}{\textbf{CLOSED}} & O30 has a finite 1.137-fold crossover; full duty has none at finite flux because the bulk pinch is zero \\
Symmetric-current benefit in the present model & \textcolor{Green}{\textbf{CLOSED}} & No useful mean EO increment; defer transient AC until measured dynamics justify it \\
Exploratory current density and area & \textcolor{Green}{\textbf{CLOSED}} & Report solved values; apply neither an A/m$^2$ ceiling nor an area acceptance gate now \\
Start material search and coupon work & \textcolor{Green}{\textbf{GO}} & Begin now using the joint-property and same-sample requirements \\
Rank a material candidate & \textcolor{Orange}{\textbf{OPEN}} & No candidate yet has the complete same-sample property set \\
Full post-ERV duty on one common bus below 2~V & \textcolor{Orange}{\textbf{OPEN}} & No confirmed sampled point: 21 physical no-solution and 28 numerical/local unresolved refined coordinates; test staging/zoning only after material data \\
Select post-ERV versus direct-ERV architecture & \textcolor{Orange}{\textbf{OPEN}} & Requires measured film, receiver, thermal, pressure-drop and durability evidence \\
Claim net system or liquid-product performance & \textcolor{Red}{\textbf{NO-GO}} & Requires a complete measured receiver/heat/hardware ledger and long-term operation \\
\bottomrule
\end{tabularx}
\end{center}

\subsection{Immediate experimental handoff}

Use the 25-cm$^2$ controlled-flow coupon as the primary transport experiment.
The current practical basis is a roughly 110$\times$110-mm fixture,
2-mm flow gaps and 6~L/min per side, followed by 3/6/12-L/min and
1/2/4-mm gap repeats. Run zero-current controls and an ABBA polarity sequence;
measure both source loss and receiver gain, charge and voltage, and keep the
membrane differential pressure within $\pm5$~Pa. Double seals, foil blanks and
aperture/perimeter masks are required to separate active-area transport from
bypass and edge leakage.

For each coupon, freeze candidate ID, batch, ion form, conditioning, temperature,
humidity and wet thickness. Measure sorption, $D$/mobility, drag, wet
conductivity, compressed ASR, film sensitivity, swelling/creep and the coupled
two-sided water balance on that same material. A quiescent closed cell may be
used only as a transient diagnostic; the present model does not support a
fixed-boundary quiescent steady-state rate.

\begin{mainbox}
\textbf{Stop criterion reached.} No further broad model study is required
before materials are sought. Continue modelling only as a replay of measured
candidate data, a pre-registered design of the coupon experiment, or closure
of a named subsystem such as the receiver/thermal ledger.
\end{mainbox}

% -----------------------------------------------------------------------------
\section{Claim boundaries and canonical evidence}
\eyebrow{What this report does and does not establish}

All active results are steady-state and conditional. The direct-ERV model
prescribes sensible effectiveness and external EO-heat rejection and does not
resolve conjugate temperature/humidity feedback, condensation, manifold
pressure drop, leakage, electrode kinetics, ageing or control dynamics. The
signed-current replay is quasi-static; it does not predict frequency response,
hysteresis, induced-charge EO, electrode rectification or transport in a
spatially graded membrane. No full-coupled or commercial claim is enabled.

\begin{stopbox}
Permitted wording is: \emph{the model identifies a conditional energy window
and a measurable membrane-property region}. Do not state that 162.4~kW or
435.9~kW has been saved, that the 80.16\% direct-ERV window is a net system
saving, that full duty is impossible, or that a candidate membrane has been
identified.
\end{stopbox}

\subsection*{Canonical evidence paths}

\begin{itemize}
  \item \evidencepath{outputs/fom_assessment_v1/system_fom.csv} and
    \evidencepath{device_fom.csv}: frozen equal-service and known-core energy/FOM ledgers;
  \item \evidencepath{outputs/workshop_common_voltage_diffusivity_v1/}:
    post-ERV common-voltage diffusivity replay below 2~V;
  \item \evidencepath{outputs/direct_active_eo_erv_study_v1/study/}:
    direct active-ERV cases, thermodynamic/engineering screens and convergence;
  \item \evidencepath{outputs/direct_erv_gradient_closure_v1/}:
    film identity, area/film sweep, signed $+J/0/-J$ replay, figures, audit and source manifest;
  \item \evidencepath{outputs/project_closure_check_v1/}:
    project-level thermodynamic reclassification, common-voltage replay and decision gates;
  \item \evidencepath{outputs/eo_vapor_cell_sizing_v1/} and
    \evidencepath{outputs/eo_vapor_boundary_analysis_v1/}:
    coupon sizing, detectability and quiescent/transient boundary analysis;
  \item \evidencepath{reports/workshop_pack_2026-08-30/MATERIAL_SEARCH_BRIEF.md}:
    material-property schema, non-gates and supplier/coupon screening sequence.
\end{itemize}

Every numerical output directory above contains an audit and/or source
manifest. Audit PASS means that the declared numerical and identity checks
passed; it does not mean that a material, module or technology passed
validation.

\vfill
{\sffamily\scriptsize\color{Slate}
End of model-closure report. This document freezes the current interpretation
of the declared model evidence as of 31 August 2026.}

\end{document}
